\section*{Capítulo \ref{ch:b3:deformation-mechanisms} --- Mecanismos de deformación y zonas de cizalla}

\begin{solution}{exo:b3:deformation-mechanisms:1}
Forsterita: $1573/2163 = 0.73$; hielo: $263/273 = 0.96$. El hielo está mucho más cerca de
su punto de fusión y se deforma por fluencia con más facilidad: los glaciares
fluyen en la superficie de la Tierra.
\end{solution}

\begin{solution}{exo:b3:deformation-mechanisms:2}
$\gamma = \tan\ang{70} = 2.75$ y $\tan\ang{80} = 5.67$. Cerca de \ang{90} la
tangente cambia muy deprisa con el ángulo: un grado de error en $\psi =
\ang{84}$ modifica $\gamma$ en una quinta parte.
\end{solution}

\begin{solution}{exo:b3:deformation-mechanisms:3}
$\sigma = (10^{3.56}/10)^{1/1.26} = 363.1^{0.794} = \qty{108}{MPa}$.
\end{solution}

\begin{solution}{exo:b3:deformation-mechanisms:4}
$E/RT = 223000/(8.314\times773.15) = 34.69$, así que $\dot\varepsilon =
\num{1.1e-4}\times20^4\times e^{-34.69} = \num{1.1e-4}\times\num{1.6e5}\times
\num{8.6e-16} = \qty{1.5e-14}{s^{-1}}$.
\end{solution}

\begin{solution}{exo:b3:deformation-mechanisms:5}
$\eta = \sigma/(2\dot\varepsilon)$: $\num{65.4e6}/\num{2e-14} =
\qty{3.3e21}{Pa.s}$ a \qty{400}{\celsius} y $\num{18.0e6}/\num{2e-14} =
\qty{9.0e20}{Pa.s}$ a \qty{500}{\celsius}: la corteza media a estas
temperaturas es aproximadamente tan viscosa como el manto bajo los casquetes de hielo.
\end{solution}

\begin{solution}{exo:b3:deformation-mechanisms:6}
$n = \log(\num{1.6e-4}/\num{e-5})/\log2 = \log16/\log2 = 4$: \omterm{def:b3:deformation-mechanisms:creep}{fluencia por dislocaciones},
cuyo \omterm{def:b3:deformation-mechanisms:flow-law}{exponente de esfuerzo} es muy superior a 1.
\end{solution}

\begin{solution}{exo:b3:deformation-mechanisms:7}
$\ln10 = (E/R)(1/1000 - 1/1100)$, luego $E = 8.314\times2.3026/\num{9.091e-5} =
\qty{211}{kJ.mol^{-1}}$.
\end{solution}

\begin{solution}{exo:b3:deformation-mechanisms:8}
Nabarro--Herring ($d^{-2}$): 100 veces más rápida; Coble ($d^{-3}$): 1000 veces.
A lo largo de los bordes, los átomos que difunden usan una trayectoria cuya sección transversal, $\delta
d$, que disminuye con el grano, además de la distancia más corta y del
menor volumen que desplazar.
\end{solution}

\begin{solution}{exo:b3:deformation-mechanisms:9}
$\sigma_b^{2.5} = \num{3.59e-6}\,d^{-3}/\num{6.54e-17}$: para $d = 1000$,
$\sigma_b = 54.9^{0.4} = \qty{5.0}{MPa}$; para $d = 100$, $\sigma_b =
\num{54900}^{0.4} = \qty{79}{MPa}$. La pendiente es $-3/2.5 = -1.2$: un
grano diez veces más fino eleva el esfuerzo de frontera en un factor de $10^{1.2} = 16$.
\end{solution}

\begin{solution}{exo:b3:deformation-mechanisms:10}
En ambos caminos la diferencia de concentración impulsora es la misma y la
distancia es $d$; el volumen que desplazar por unidad de deformación es $d^3$. La red cristalina
ofrece una sección transversal $\propto d^2$, así que la velocidad va como $d^2/(d\,d^3) =
d^{-2}$; los bordes ofrecen $\propto\delta d$, así que $\delta d/(d\,d^3) = \delta
d^{-3}$. La razón entre las velocidades de borde y de red es $(B_B\delta D_B)/(B_LD_Ld)$,
mayor que 1 cuando $d < (B_B/B_L)\,\delta D_B/D_L$.
\end{solution}

\begin{solution}{exo:b3:deformation-mechanisms:11}
$\eta \propto \dot\varepsilon^{(1-n)/n} = \dot\varepsilon^{-2.5/3.5}$: diez veces
más deprisa divide la viscosidad por $10^{0.714} = 5.2$. Una banda que casualmente
se deforma más deprisa se vuelve más débil, así que absorbe más de la deformación:
la deformación se concentra en ella.
\end{solution}

\begin{solution}{exo:b3:deformation-mechanisms:12}
Antes: $\dot\varepsilon_{\mathrm{dis}} = \num{6.54e-17}\times10^{3.5} =
\qty{2.07e-13}{s^{-1}}$ y $\dot\varepsilon_{\mathrm{dif}} = \num{3.59e-6}
\times10/3000^3 = \qty{1.3e-15}{s^{-1}}$: \qty{2.08e-13}{s^{-1}}, \omterm{def:b3:deformation-mechanisms:creep}{fluencia por
dislocaciones}. Después: $\dot\varepsilon_{\mathrm{dif}} = \num{3.59e-6}\times10/30^3 =
\qty{1.33e-9}{s^{-1}}$, \omterm{def:b3:deformation-mechanisms:creep}{fluencia por difusión}, \num{6400} veces más rápida con el mismo
esfuerzo: la recristalización ha debilitado enormemente la zona.
\end{solution}

\begin{solution}{pb:b3:deformation-mechanisms:1}
\textbf{1.} El dique gira desde \ang{90} hasta \ang{6} respecto de la zona: $\psi =
\ang{84}$ y $\gamma = \tan\ang{84} = 9.51$.

\textbf{2.} $\gamma w = 9.51\times400 = \qty{3804}{m}$, unos \qty{3.8}{km}.

\textbf{3.} $S_1 = \sqrt{1 + 9.51^2/4} + 9.51/2 = 4.859 + 4.755 = 9.614$,
$S_2 = 1/S_1 = 0.104$, y $S_1/S_2 = 92.4$.

\textbf{4.} $\tan2\theta' = 2/9.51 = 0.210$, $\theta' = \ang{5.9}$: la
foliación, próxima al eje largo, gira cada vez más hacia la zona a medida que $\gamma$
crece.

\textbf{5.} Siguiendo el dique hacia el norte, este es arrastrado hacia el este dentro de la zona y
sale de ella más al este: la pared norte se ha desplazado hacia el este respecto de la sur.
Visto desde la pared sur, el lado opuesto se desplazó hacia la derecha: la zona es
dextral.

\textbf{6.} En el plano que contiene la lineación de estiramiento, horizontal aquí
por tratarse de una zona de desgarre, y perpendicular a la foliación vertical: una
sección horizontal. Su foliación S se inclina en sentido antihorario respecto a la zona, hacia
el este-noreste, como el eje largo de la elipse de deformación.

\textbf{7.} $\sigma = (10^{3.56}/20)^{1/1.26} = 181.5^{0.794} = \qty{62}{MPa}$.

\textbf{8.} $242.1^{0.794} = \qty{78}{MPa}$ para \qty{15}{\micro m} y
$145.2^{0.794} = \qty{52}{MPa}$ para \qty{25}{\micro m}.

\textbf{9.} Sí: los experimentos abarcaron entre 34 y \qty{268}{MPa}.

\textbf{10.} La zona se deformó a \qty{450}{\celsius}, muy por debajo de los 800 a
\qty{1100}{\celsius} de los experimentos; una relación que no depende de
la temperatura puede trasladarse hasta ella, aunque sigue siendo una
extrapolación.

\textbf{11.} $D = 3631/100^{1.26} = 3631/331 = \qty{11}{\micro m}$.

\textbf{12.} $E/RT = 223000/(8.314\times723.15) = 37.09$ y $e^{-37.09} =
\num{7.79e-17}$.

\textbf{13.} $\dot\varepsilon = \num{1.1e-4}\times62^4\times\num{7.79e-17} =
\num{1.1e-4}\times\num{1.478e7}\times\num{7.79e-17} = \qty{1.27e-13}{s^{-1}}$.

\textbf{14.} $\dot\gamma = 2\dot\varepsilon = \qty{2.53e-13}{s^{-1}}$.

\textbf{15.} $\eta = \num{62e6}/(2\times\num{1.27e-13}) = \qty{2.4e20}{Pa.s}$,
próxima a los \qty{3e20}{Pa.s} del manto superior bajo el lago Ardo: la zona
era tan débil como el manto.

\textbf{16.} $e^{-E/RT}$ cambia en un factor de 0.064 a \qty{400}{\celsius}
(dieciséis veces más lento) y 11.0 a \qty{500}{\celsius}.

\textbf{17.} $(78/62)^4 = 2.5$ veces más rápido.

\textbf{18.} $t = 9.51/\num{2.53e-13} = \qty{3.75e13}{s}$.

\textbf{19.} $\num{3.75e13}/\num{3.156e10} = \qty{1190}{kyr}$.

\textbf{20.} $\qty{3804}{m}/\qty{1.19e6}{yr} = \qty{3.2}{mm/yr}$.

\textbf{21.} Diez o más veces más lenta que los principales límites de placa: la zona
era un ramal secundario, o uno de varios que compartían el movimiento de un límite.

\textbf{22.} $\dot\gamma = 0.02/(\num{3.156e7}\times400) = \qty{1.58e-12}{s^{-1}}$,
6.2 veces más rápido; el esfuerzo aumenta en $6.2^{1/4} = 1.58$, hasta \qty{98}{MPa}, y
los granos tendrían un tamaño de $3631/98^{1.26} = \qty{11}{\micro m}$.

\textbf{23.} El tiempo varía como $\sigma^{-4}$: el intervalo de tamaños de grano por sí solo da
entre $(62/78)^4 = 0.40$ y $(62/52)^4 = 2.0$ veces la estimación. A ello se suma la
incertidumbre del exponente $1.26 \pm 0.13$ y de las constantes de la \omterm{def:b3:deformation-mechanisms:flow-law}{ley de flujo}, que
la amplía aún más: el resultado es un orden de magnitud, no una fecha.

\textbf{24.} La \omterm{def:b3:deformation-mechanisms:shear-zone}{zona de cizalla} de Varden tardó unos
\textbf{\qty[text-series-to-math=true]{1190}{kyr}} en acumular su deformación.
\end{solution}
